\section{Related work}
\label{sec:related}

\paragraph{Haze forecasting in Equatorial Asia.}
Regional haze warnings rely on dispersion models and on satellite hotspot and smoke-plume analyses
from the ASEAN Specialised Meteorological Centre. Process studies link the transport of Indonesian
smoke to Singapore to south-easterly low-level jets and boundary-layer mixing
\citep{huang2024haze}. On the data-driven side, \citet{wang2020haze} trained a convolutional network
on 35 years of meteorological and hydrological fields to forecast severe haze at Singapore. They
found skill above climatology but many missed events, and they used no fire data. Fire-to-PM
modelling in Riau and Pontianak with gradient boosting and ensembles exists \citep{pontianak2026},
but it targets PM\textsubscript{2.5} at a single city and does not isolate directional fire
information.

\paragraph{Visibility as an observed target.}
Airport visibility has long been used as a proxy for particulate pollution where monitors are
scarce, including reconstructions of PM\textsubscript{2.5} from near-surface visibility across
the Northern Hemisphere since 1959 \citep{metarpm2024}. In Central Kalimantan, visibility at Palangka Raya falls non-linearly with
PM\textsubscript{2.5} and PM\textsubscript{10}, dropping below 0.2 miles during the 2015 fires
\citep{palangka2022}. Tree ensembles also forecast airport fog and low visibility from METAR and
reanalysis \citep{fogtransfer2025}. We use visibility with haze codes as the forecast target
itself.

\paragraph{Wind-aware fire exposure.}
Weighting fire radiative power by distance and by the angle between wind and fire-to-receptor
direction is an established feature-engineering choice. \citet{gnnfire2023} use inverse-distance
and wind-angle weighting in graph-network PM\textsubscript{2.5} forecasts, and \citet{sci2026exposure}
evaluate distance- and direction-aware exposure metrics. Epidemiological work defines upwind
exposure from wind direction and fire location \citep{aguilera2021}. In contrast, a recent
site-level forecasting study found that inverse-distance fire features added nothing beyond
prior-day PM\textsubscript{2.5} and weather \citep{santacruz2026}. To our knowledge, no study has
isolated the \emph{directional} term against a non-directional control with identical weighting,
checked it with rotated-wind placebos, and repeated the test across several haze seasons and
fire-data latencies. That is the gap this paper fills.
